\documentclass[12pt,a4paper]{article}
\usepackage[utf8]{inputenc}
\usepackage[T1]{fontenc}
\usepackage[brazil]{babel}
\usepackage{amsmath, amssymb}
\usepackage{geometry}
\geometry{top=2cm, bottom=2.5cm, left=2cm, right=2cm}
\usepackage{enumitem}
\setlist{nosep}

\title{Exercícios de Análise Combinatória \\ com Explicações Detalhadas}
\author{}
\date{}

\begin{document}
\maketitle

\section*{1. Números de 3 algarismos distintos com os dígitos 1, 2, 3, 4, 5}

\textbf{Raciocínio:}  
Queremos formar números de 3 dígitos distintos.  
- O primeiro dígito (centena) pode ser qualquer um dos 5 disponíveis.  
- O segundo dígito (dezena) deve ser diferente do primeiro, sobrando 4 opções.  
- O terceiro dígito (unidade) deve ser diferente dos dois anteriores, sobrando 3 opções.  

\[
\text{Total de números} = 5 \times 4 \times 3 = 60
\]

\textbf{Resposta:} 60 números

\section*{2. Escolha de 5 questões entre 8}

\textbf{Raciocínio:}  
Aqui a ordem não importa, apenas a escolha. Podemos pensar assim: escolher 5 questões é o mesmo que escolher 3 que não serão respondidas.  

\[
\text{Total de maneiras} = \frac{8 \times 7 \times 6}{3 \times 2 \times 1} = 56
\]

\textbf{Resposta:} 56 maneiras

\section*{3. Anagramas da palavra ``MATEMÁTICA''}

\textbf{Raciocínio:}  
A palavra tem 10 letras, mas algumas se repetem: A aparece 3 vezes, M aparece 2 vezes, T aparece 2 vezes.  
Quando há repetição, dividimos o fatorial pelo fatorial de cada letra repetida para evitar contar duplicações.  

\[
\text{Total de anagramas} = \frac{10!}{3! \cdot 2! \cdot 2!} = 151\,200
\]

\textbf{Resposta:} 151\,200 anagramas

\section*{4. Comissões em sala com 10 homens e 8 mulheres}

\textbf{(a) Qualquer pessoa:}  
Escolhemos 4 pessoas entre 18 (10 homens + 8 mulheres).  

\[
\text{Total} = \frac{18 \times 17 \times 16 \times 15}{4 \times 3 \times 2 \times 1} = 3\,060
\]

\textbf{(b) 2 homens e 2 mulheres:}  
- Homens: escolher 2 de 10: \(\frac{10 \times 9}{2} = 45\)  
- Mulheres: escolher 2 de 8: \(\frac{8 \times 7}{2} = 28\)  
Multiplicamos os dois resultados: \(45 \times 28 = 1\,260\)

\textbf{(c) Pelo menos 1 mulher:}  
Podemos usar o princípio complementar: total de comissões menos comissões sem mulheres.  
- Comissões só de homens: \(\frac{10 \times 9 \times 8 \times 7}{4 \times 3 \times 2 \times 1} = 210\)  
\(\text{Com pelo menos 1 mulher} = 3\,060 - 210 = 2\,850\)

\section*{5. Números pares de 4 algarismos distintos com dígitos 1--6}

\textbf{Raciocínio:}  
- Último dígito deve ser par (2, 4, 6): 3 opções  
- Primeiro dígito: qualquer dígito restante: 5 opções  
- Segundo dígito: qualquer dígito restante: 4 opções  
- Terceiro dígito: 3 opções restantes  

\[
\text{Total} = 3 \times 5 \times 4 \times 3 = 180
\]

\textbf{Resposta:} 180 números pares

\section*{6. Restaurante: escolha de refeição}

\textbf{(a) Um item de cada tipo:}  
\[
5 \text{ saladas} \times 3 \text{ pratos principais} \times 4 \text{ sobremesas} = 60
\]

\textbf{(b) Apenas prato principal e sobremesa:}  
\[
3 \times 4 = 12
\]

\section*{7. Diagonais de um decágono}

\textbf{Raciocínio:}  
- Um décagono tem 10 vértices.  
- Cada vértice se conecta a outros 7 (excluindo ele mesmo e os dois vizinhos).  
- Dividimos por 2 para não contar duas vezes a mesma diagonal.  

\[
\text{Total} = \frac{10 \times 7}{2} = 35
\]

\section*{8. Pessoas em mesa circular (6 pessoas)}

\textbf{Raciocínio:}  
- Em mesas circulares, fixamos uma pessoa como referência, restando \(5!\) formas de organizar as demais.  

\[
5! = 120
\]

\section*{9. Senhas com 2 letras e 2 números}

\textbf{(a) Letras e números podem repetir:}  
\[
26^2 \times 10^2 = 67\,600
\]

\textbf{(b) Nenhum caractere pode repetir:}  
\[
26 \times 25 \times 10 \times 9 = 58\,500
\]

\section*{10. Subconjuntos de um conjunto com 5 elementos}

\textbf{Raciocínio:}  
- Cada elemento pode estar ou não no subconjunto: 2 opções  
\[
2^5 = 32
\]

\section*{11. Números ímpares de 4 dígitos distintos entre 1000 e 5000}

\textbf{Raciocínio:}  
Dividimos em casos pelo primeiro dígito (milhar):  
- Milhar 1 ou 3: unidade ímpar ≠ milhar (4 opções), duas casas do meio: 8 × 7  
- Milhar 2 ou 4: unidade ímpar (5 opções), duas casas do meio: 8 × 7  

\[
\text{Total} = 448 + 560 = 1\,008
\]

\section*{12. Comissão de 5 com maioria feminina (7 homens, 5 mulheres)}

\begin{itemize}
    \item 3 mulheres + 2 homens: \( \frac{5 \times 4 \times 3}{3!} \times \frac{7 \times 6}{2} = 210 \)  
    \item 4 mulheres + 1 homem: 35  
    \item 5 mulheres: 1
\end{itemize}

\[
\text{Total} = 246
\]

\section*{13. Anagramas de ``PARALELOGRAMO'' que começam e terminam com vogal}

\textbf{Raciocínio:}  
Escolhemos vogal inicial e final, depois permutamos as letras do meio considerando repetições.  

\textbf{Resultado aproximado:} 24\,948\,000 anagramas

\section*{14. Jogos em torneio de 8 times}

\[
\frac{8 \times 7}{2} = 28 \text{ jogos}
\]

\section*{15. Números de 1 a 1000 divisíveis por 3 ou 5}

\[
\text{Divisíveis por 3} = 333, \quad \text{por 5} = 200, \quad \text{por 15} = 66
\]  
\[
\text{Total} = 333 + 200 - 66 = 467
\]

\section*{16. Distribuição de 8 livros entre 3 estudantes, cada um recebe pelo menos 2 livros}

\begin{itemize}
    \item Caso (2,2,4): 1\,260 formas  
    \item Caso (2,3,3): 1\,680 formas
\end{itemize}

\[
\text{Total} = 2\,940
\]

\section*{17. Triângulos em octógono}

\[
\frac{8 \times 7 \times 6}{3 \times 2 \times 1} = 56
\]

\section*{18. Soluções naturais da equação x+y+z=12}

\begin{itemize}
    \item Naturais $\ge 1$: 55 soluções  
    \item Naturais $\ge 0$: 91 soluções
\end{itemize}

\section*{19. Organização de livros por matéria (5 de Matemática, 4 de Física, 3 de Química)}

\textbf{Raciocínio:}  
- Livros de mesma matéria juntos, tratamos cada grupo como um bloco:  
- Permutação dos blocos: 3!  
- Permutação interna dos livros: 5! × 4! × 3!  

\[
3! \times 5! \times 4! \times 3! = 103\,680
\]

\section*{20. Números de 5 dígitos com pelo menos um dígito par}

\textbf{Raciocínio:}  
- Total de números de 5 dígitos: 9 × 10^4 = 90\,000  
- Total apenas ímpares: 5^5 = 3\,125  
- Pelo menos 1 dígito par: \(90\,000 - 3\,125 = 86\,875\)

\end{document}

